\scanpage{27}
\mspage{1}

\uline{Structures stratifiées.}

1. \uline{La situation la plus élémentaire}, en un sens qui apparaîtra, sera la suivante. On se donne des espaces (disons !) en attendant que \unsure{ça|ce} devienne des topos\ldots)

(1)
\diagram{\begin{tikzcd}[ampersand replacement=\&]Y\arrow[r,hook]\&X\end{tikzcd}}

$Y$ fermé dans $X$. On suppose que la ``structure normale\unsure{''} de $X$ le long de $Y$ soit loc. triviale''\unsure{)}, i.e. que pour tout $y\in Y$, existe un voisinage ouvert $X'$ de $y$ \unsure{et} un espace $Z$ pointé par $z$, et un diagramme commutatif (où $Y'=Y\cap X'$)

(2)
\diagram{\begin{tikzcd}[ampersand replacement=\&,column sep=large,row sep=large]
Y'\arrow[r,hook]\arrow[d,"\text{can}"',"\sim"]\&X'\arrow[d,"\sim"]\\
Y'\times\{z\}\arrow[r,hook,"\mathrm{id}_{Y'}\times i_{z,Z}"']\&Y'\times Z
\end{tikzcd}}

où $X'\simeq Y'\times Z$ est un homéomorphisme.

On considère alors ``le voisinage tubulaire'' (``infinitésimal'') de $Y$ dans $X$, soit

(3)
\diagram{\begin{tikzcd}[ampersand replacement=\&]
Y\arrow[r,hook]\&\unsure{\mathcal V_{Y,X}|V_{Y,X}}\arrow[r,hook]\&X
\end{tikzcd}},

dont il faudra donner une définition

\margin{26}
